\begin{textsample}{POS Dim 2 – go – Score 26.00 – go/go\_em\_17.txt}  \label{ex:f2_pos_175}
3.3 - Considerações Acerca dos Aspectos e das Condições Necessárias à Implementação dos \textbf{Itinerários} de Formação \textbf{Técnica} \textbf{Profissional} Ao se pensar o planejamento e a implementação de \textbf{itinerários} \textbf{formativos} no campo da Formação \textbf{Técnica} e Profissional, é importante considerar e refletir acerca dos aspectos abordados na sequência deste texto, os quais constituem as condições básicas para a \textbf{efetivação} das propostas de formação desta natureza.
Esses devem ser tomados, ao mesmo tempo, como desafios a serem enfrentados, mas também como elementos \textbf{orientadores} deste processo de preparação das propostas de \textbf{itinerários} \textbf{formativos} no âmbito da Educação Profissional.
Tais aspectos, em sua maioria, já constituem objetos de preocupação do dispositivo responsável por estabelecer as \textbf{diretrizes} curriculares para o Ensino Médio, conforme Resolução CNE/CEB nº 3/18, art. 12, Inciso V, § 1º : > Os \textbf{itinerários} \textbf{formativos} devem considerar as demandas e necessidades do mundo contemporâneo, estar \textbf{sintonizados} com os diferentes interesses dos estudantes e sua inserção na sociedade, o contexto local e as possibilidades de \textbf{oferta} dos sistemas e instituições de ensino ( BRASIL, 2018 ).
Os aspectos a serem considerados são : - Concepção/elaboração das propostas pedagógicas - planos de \textbf{curso} ; - Sintonia com o contexto socioprodutivo local e regional ; - Infraestrutura física e de equipamentos das instituições ; - Material instrucional específico ; - Recursos humanos minimamente qualificados.
Planejamento e Elaboração das Propostas Pedagógicas ( Planos de \textbf{Curso} ) Este aspecto abrange todos os demais aspectos na sequência desta discussão, pois, na elaboração de uma proposta pedagógica de \textbf{curso}, há que se considerar, além dos fundamentos e princípios teórico-metodológicos, os aspectos de caráter prático e operacional sobre os quais aportam as ações pedagógicas e que possuem enorme peso no que se refere à viabilidade da implementação da proposta pretendida.
O ponto de partida do planejamento diz respeito à definição de qual \textbf{curso} \textbf{ofertar} e o tipo, bem como à estratégia de \textbf{oferta}, devendo considerar o potencial e as \textbf{vocações} socioeconômicas locais e regionais, a capacidade da instituição de realizar a \textbf{oferta} e os interesses dos/as estudantes, público-alvo da proposta.
Feito isto, o passo seguinte é a elaboração da proposta pedagógica de \textbf{curso}, neste caso, concebida na perspectiva e como estratégia de \textbf{itinerário} \textbf{formativo} do Ensino Médio, a qual deve começar pela definição do \textbf{perfil} \textbf{Profissional} de conclusão dos/as egressos/as.
O \textbf{perfil} \textbf{Profissional} de conclusão constitui o ponto inicial e de referência para todo o planejamento da proposta e sua definição.
Como citado anteriormente, poderá ser orientado tanto pelos Catálogos/Guias Nacionais de \textbf{Cursos} e/ou pela própria CBO quanto por levantamentos específicos junto aos próprios setores produtivos acerca das competências e habilidades que tais agentes produtivos esperam encontrar no \textbf{perfil} de conclusão dos/as egressos/as do \textbf{curso}.
Uma vez definido o \textbf{perfil} de conclusão, as demais etapas do planejamento da proposta, como justificativa, objetivos, a organização curricular e as respectivas metodologias de ensino, dentre outras, serão desenvolvidas com base no \textbf{perfil} de conclusão, à luz dos princípios norteadores da EPT e dos eixos \textbf{estruturantes} dos \textbf{itinerários} \textbf{formativos}, sobre os quais trataremos mais adiante.
Os princípios norteadores da EPT, dispostos na Resolução CNE/CEB nº 6/2012, dispositivo este que estabelece as \textbf{diretrizes} curriculares para a Educação Profissional de Nível \textbf{Técnico}, podem ser tomados como balizadores de toda esta modalidade educacional.
Dentre eles destacam -se : 1.
A articulação da EPT com os demais níveis e modalidade de ensino - a educação \textbf{Profissional}, conforme atualmente concebida, não visa uma formação assistencialista, nem limitada à formação \textbf{técnica} para atendimento do mercado de trabalho.
Pretende, isto sim, em articulação com os outros segmentos da educação e com as dimensões do trabalho, da ciência e da tecnologia, promover uma formação ampla, integradora e contextualizada que, além dos saberes, competências e habilidades específicos de uma dada profissão, preocupa -se, também, com o desenvolvimento e a promoção de princípios e valores necessários à compreensão global do processo produtivo e ao exercício da cidadania em uma sociedade que se pretende : democrática, solidária e inclusiva ; 2.
O respeito aos valores estéticos, políticos e éticos da educação nacional - os valores estéticos, expressos na estética da sensibilidade, estão diretamente relacionados aos conceitos de \textbf{qualidade} e respeito aos/às outros/as.
O respeito pelo/a outro/a exige o desenvolvimento de uma cultura do trabalho centrada no gosto pelo trabalho bem feito e acabado, quer na prestação de serviços, quer na produção de bens ou de conhecimentos, não transigindo com o trabalho mal feito e inacabado ( Parecer do CNE 16/99, p. 292 ), mas que prima pela \textbf{qualidade} daquilo que é feito, produzido, elaborado.
Os valores políticos, expressos na \textbf{política} da igualdade, numa síntese breve, são referentes ao direito de todos/as à educação para o trabalho, numa perspectiva de superação da dicotomia : trabalho intelectual ( criação ), trabalho manual ( execução ) e de formação de trabalhadores/as aptos/as à tomada de decisão e à intervenção nos processos de trabalho, sendo esta, inclusive, uma exigência fundamental do novo contexto do mundo do trabalho que tem passado por profundas alterações em função das constantes inovações tecnológicas e das atuais formas de reorganização dos processos de produção e do trabalho.
Os valores éticos, expressos na ética da identidade de valores relacionados à \textbf{corresponsabilidade} e \textbf{observância} de suas atribuições conforme estabelecido entre as partes.
Em termos \textbf{Profissionais} e do trabalho, significa que se deve respeitar as normas e regras sócio organizacionais, compreendendo que os direitos \textbf{Profissionais} devem advir em função da competência do mérito, sem qualquer tipo de favoritismo ou privilégio. 3.
O trabalho assumido como princípio educativo - este princípio parte do pressuposto de que o ser humano se educa na relação com outros seres humanos e com a natureza à medida que precisa produzir as condições necessárias à manutenção da vida individual e social.
No campo da ciência da educação, segundo Kuenzer ( 1989 ), o trabalho como princípio educativo consiste em uma diretriz mais geral de uma proposta de educação que reunifique cultura e produção/educação e trabalho.
O trabalho, nesta perspectiva, deve ser tomado em seu sentido ontológico, ou seja, como atividade inerente ao próprio ser humano.
Conforme a autora, uma escola fundada neste princípio, teria as seguintes características \textbf{estruturantes}.
Do ponto de vista : a ) da estrutura, ela seria única ; b ) do conteúdo, ela seria politécnica ; c ) do método, ela seria teórico-prática ( dialética ) ; d ) da gestão, ela seria democrática ; e ) das condições físicas, ela seria moderna e \textbf{atualizada}, com equipamentos, laboratórios e bibliotecas que permitiriam a apropriação do saber científico, tecnológico e histórico-crítico da sociedade moderna. 4.
A pesquisa como princípio pedagógico - consiste em estratégia metodológica que faz com que as aprendizagens ocorram de forma mais criativa, autônoma e significativa, tornando -se mais efetivas.
A problematização da realidade, o exercício da investigação científica acerca dos problemas formulados e as respostas suscitadas ao final do processo de análise dos dados encontrados, conduzem à construção de novos conhecimentos e à verdadeira aprendizagem e, ao que é mais importante, possibilitar que o/a estudante emerja como protagonista deste processo ; 5.
A indissociabilidade entre teoria e prática no processo de ensino-aprendizagem - teoria e prática são dimensões distintas de um mesmo processo em que os saberes, o desenvolvimento de competências e as habilidades práticas ( saber fazer ) requeridos pela natureza do trabalho, constituem um movimento único e interdependente que, assim assumidas, possibilitam a realização do propósito de produzir aprendizagens efetivas e significativas.
Consiste em estratégia metodológica baseada na retroalimentação entre o saber teórico e o saber fazer ( prática ) que facilita e propicia aos/às estudantes uma melhor compreensão dos objetos de estudos propostos ; 6.
A contextualização, flexibilidade e interdisciplinaridade - são princípios que, aplicados às estratégias educacionais, favorecem a compreensão do mundo físico e sociocultural contemporâneo, por sinal, em constante processo de mudança.
Assim, considerando uma realidade em constante processo de transformação, é extremamente necessário a atualização dos \textbf{currículos} e a renovação e diversificação das metodologias de ensino-aprendizagem, bem como a interação/articulação entre os diferentes campos de conhecimento que, em última instância, constituem partes de um todo significante ; 7.
A articulação com o desenvolvimento socioeconômico e ambiental da região - em se tratando da Educação Profissional, é imprescindível a contextualização socioprodutiva local e regional de onde os \textbf{cursos} pretendem ser \textbf{ofertados}, de modo que haja consonância entre seus \textbf{perfis} de conclusão e a realidade do cenário social, produtivo e ambiental, a fim de que os/as egressos/as dos \textbf{cursos} possam encontrar colocação no mercado de trabalho local e regional e exercer suas profissões na perspectiva da responsabilidade social e ambiental requeridas pela lógica do desenvolvimento sustentável ; 8.
O reconhecimento das diversidades socioculturais - ( respeito às formas de ser e pensar das pessoas e grupos étnico-raciais ) e produtivas ( diferentes às formas de organização dos processos produtivos e do trabalho ) ; 9.
A autonomia da instituição educacional na concepção, elaboração, execução, avaliação e revisão do seu Projeto Político-Pedagógico ( PPP ) - o PPP deve ser entendido como uma construção do coletivo escolar, à luz dos dispositivos legais e das \textbf{diretrizes} curriculares pertinentes.
Nesse aspecto, os planos de \textbf{curso} a serem elaborados precisam, necessariamente, estar em consonância com o PPP da instituição, fazendo parte de sua intencionalidade educativa e de sua identidade institucional.
Assim, as propostas pedagógicas para os \textbf{itinerários} da Formação \textbf{Técnica} e Profissional precisam ser elaboradas à luz dos princípios supracitados, contar com a participação de \textbf{Profissionais} de área dos respectivos \textbf{cursos} ora definidos e com o envolvimento de toda a equipe da unidade de ensino, para que possam \textbf{atender} às expectativas da formação desejada para o Ensino Médio, conforme a BNCC ( BRASIL, 2018, p. 479 ) : > O conjunto dessas aprendizagens ( Formação Geral Básica e Itinerário \textbf{Formativo} ) deve \textbf{atender} às finalidades do Ensino Médio e às demandas de \textbf{qualidade} de formação na contemporaneidade, bem como às expectativas presentes e futuras das juventudes.
Por fim, há que se atentar para o aspecto legal e de legitimação das propostas concebidas, uma vez que precisam ser submetidas aos \textbf{órgãos} competentes, conforme estabelece a LDBEN nº 9.394/96 em seu artigo 36, § 8º : > A \textbf{oferta} de formação \textbf{técnica} e \textbf{Profissional} a que se refere o inciso V do caput, realizada na própria instituição ou em parceria com outras instituições, deverá ser aprovada previamente pelo Conselho Estadual de Educação, \textbf{homologada} pelo \textbf{Secretário} Estadual de Educação e certificada pelos sistemas de ensino ( BRASIL, 1996 ).
A seguir, serão tratados os demais aspectos básicos a serem considerados no processo de planejamento, elaboração e implementação das propostas educacionais referentes ao \textbf{itinerário} de Formação \textbf{Técnica} e Profissional.

% matched lemmas: atender, atualizar, corresponsabilidade, currículo, cursar, curso, diretriz, efetivação, estruturante, formativo, homologar, itinerário, observância, oferta, ofertar, orientador, perfil, política, profissional, qualidade, secretário, sintonizar, técnico, vocação, órgão
\end{textsample}
